\section*{第 \ref{ch:g7:states-of-matter} 章——物态与物态变化}

\begin{solution}{exo:g7:states-of-matter:1}
\omterm{def:g3:solids-liquids-gases:solid}{固体}：\omterm{def:g7:states-of-matter:molecule}{分子}整齐堆叠，原地\omterm{def:g3:sound-around-us:vibration}{振动}。
\omterm{def:g3:solids-liquids-gases:liquid}{液体}：\omterm{def:g7:states-of-matter:molecule}{分子}相互接触但杂乱无序，彼此滑过；
\omterm{def:g3:solids-liquids-gases:gas}{气体}：\omterm{def:g7:states-of-matter:molecule}{分子}相距很远，向各个方向快速飞动。
\end{solution}

\begin{solution}{exo:g7:states-of-matter:2}
在\omterm{def:g3:solids-liquids-gases:solid}{固体}中，每个\omterm{def:g7:states-of-matter:molecule}{分子}都守住自己在堆叠中的位置，
因此整体保持形状。在\omterm{def:g3:solids-liquids-gases:liquid}{液体}中，\omterm{def:g7:states-of-matter:molecule}{分子}仍然相互接触，
却可以自由地彼此滑过：容器中的\omterm{def:g7:states-of-matter:molecule}{分子}群下滑、填满容器，
直到再也无法下滑——形成水平的液面。
\end{solution}

\begin{solution}{exo:g7:states-of-matter:3}
\omterm{def:g2:air-around-us:air}{空气}主要是飞行的\omterm{def:g7:states-of-matter:molecule}{分子}之间的空隙：推压会把
这些飞行者挤得更近（它们的撞击会反抗——好似弹簧）。
水的\omterm{def:g7:states-of-matter:molecule}{分子}已经相互接触，没有空隙可以挤掉，
因此活塞会碰到一堵墙。
\end{solution}

\begin{solution}{exo:g7:states-of-matter:4}
\omterm{def:g3:solids-liquids-gases:gas}{气体}\omterm{def:g7:states-of-matter:molecule}{分子}相距很远，之间是大片空旷——
因此整间房的\omterm{def:g3:solids-liquids-gases:gas}{气体}都能被\omterm{prop:g7:states-of-matter:squeeze}{压缩}到一起。水的\omterm{def:g7:states-of-matter:molecule}{分子}
已经相互接触，是不可\omterm{prop:g7:states-of-matter:squeeze}{压缩}的：任何泵都无法将它压得更紧。
\end{solution}

\begin{solution}{exo:g7:states-of-matter:5}
加热使堆叠中的\omterm{def:g7:states-of-matter:molecule}{分子}\omterm{def:g3:sound-around-us:vibration}{振动}得越来越剧烈；
到达\omterm{def:g4:changes-of-state:melting-point}{熔点}时，\omterm{def:g3:sound-around-us:vibration}{振动}压倒了堆叠的束缚，
行列瓦解，\omterm{def:g7:states-of-matter:molecule}{分子}开始滑动：
\omterm{def:g3:solids-liquids-gases:solid}{固体}便变成了\omterm{def:g3:solids-liquids-gases:liquid}{液体}。
\end{solution}

\begin{solution}{exo:g7:states-of-matter:6}
物态变化只是重新排列同样的\omterm{def:g7:states-of-matter:molecule}{分子}——\omterm{def:g7:states-of-matter:molecule}{分子}既不会产生，
也不会消失，因此\omterm{def:g3:measuring-mass:mass}{质量}不会改变。要使\omterm{def:g3:measuring-mass:mass}{质量}减少，
\omterm{def:g7:states-of-matter:molecule}{分子}就必须消失（或逃出容器——
这是未盖瓶盖的瓶子造成的假象）。
\end{solution}

\begin{solution}{exo:g7:states-of-matter:7}
\omterm{def:g7:states-of-matter:molecule}{分子}运动的活跃程度。热的物体与冷的物体接触时，
活跃的\omterm{def:g7:states-of-matter:molecule}{分子}撞击迟缓的\omterm{def:g7:states-of-matter:molecule}{分子}，每次碰撞都传递一点活力，
活跃者平静下来，迟缓者随之加快——
直至两群\omterm{def:g7:states-of-matter:molecule}{分子}步调一致，\omterm{def:g3:temperature-thermometers:temperature}{温度}相等。
\end{solution}

\begin{solution}{exo:g7:states-of-matter:8}
香水\omterm{def:g7:states-of-matter:molecule}{分子}飞出，与\omterm{def:g2:air-around-us:air}{空气}\omterm{def:g7:states-of-matter:molecule}{分子}相撞、反弹、游荡，
一步步曲折地扩散到整个房间。微风会整团地裹挟它们：
顺风而行，胜过在人群中
跌跌撞撞地穿行。
\end{solution}

\begin{solution}{exo:g7:states-of-matter:9}
只有\emph{最快}的\omterm{def:g7:states-of-matter:molecule}{分子}才能挣脱\omterm{def:g3:solids-liquids-gases:liquid}{液体}的束缚，
因此\omterm{def:g2:ice-water-steam:evapcond}{蒸发}是对活跃程度的一种“征税”：逃逸者带走的
超过了它们应有的份额，留下的\omterm{def:g7:states-of-matter:molecule}{分子}群平均而言
变慢了——也就是变凉了。汗水和湿手指
正是靠带走最活跃的\omterm{def:g7:states-of-matter:molecule}{分子}而降温。
\end{solution}

\begin{solution}{exo:g7:states-of-matter:10}
阳光下，被困的\omterm{def:g7:states-of-matter:molecule}{分子}飞得更快，撞击皮肤更猛、更频繁：
如冰雹般的撞击使气球膨胀，直到绷紧的表皮与之平衡。
寒冷中，同样的\omterm{def:g7:states-of-matter:molecule}{分子}变慢，撞击减弱，
外界\omterm{def:g2:air-around-us:air}{空气}的压力便把
气球压小。
\end{solution}

\begin{solution}{exo:g7:states-of-matter:11}
水结冰时，\omterm{def:g7:states-of-matter:molecule}{分子}锁定成一种异常疏松、开放的堆叠，
空隙比流动的\omterm{def:g7:states-of-matter:molecule}{分子}群中更多，因此同样多的\omterm{def:g7:states-of-matter:molecule}{分子}
占据更大的\omterm{def:g6:mass-and-volume:volume}{体积}：\omterm{def:g6:density-floating:density}{密度}小于\omterm{def:g3:solids-liquids-gases:liquid}{液体}，所以冰会浮起。
若水是普通\omterm{def:g3:solids-liquids-gases:liquid}{液体}，冰会紧密堆叠，形成时就会\omterm{def:g1:floating-and-sinking:words}{下沉}，
池塘会从底部向上结冰——那就不会有避风的水
留给鱼了。
\end{solution}

\begin{solution}{exo:g7:states-of-matter:12}
没有任何东西搅动水，墨水却在移动：\omterm{def:g3:solids-liquids-gases:liquid}{液体}自身的\omterm{def:g7:states-of-matter:molecule}{分子}
处于永不停歇的运动之中，推挤着墨水\omterm{def:g7:states-of-matter:molecule}{分子}，
一步步随机地穿过\omterm{def:g7:states-of-matter:molecule}{分子}群——水和墨水两种\omterm{def:g3:solids-liquids-gases:liquid}{液体}
每时每刻都充满运动，即使玻璃杯看起来
完全静止。\omterm{def:g3:temperature-thermometers:temperature}{温度}升高，每次推挤都更猛，扩散也就更快。
耐心的墨水由此使看不见的舞动变得可见——
这一证据使那些怀疑\omterm{def:g7:states-of-matter:molecule}{分子}存在的人信服。
\end{solution}

\begin{solution}{pb:g7:states-of-matter:1}
\textbf{1.} 整齐排成若干行，肩并肩，每人原地
抖动——谁都不许换位置或离开自己的位置。
\textbf{2.} 队列瓦解：学生仍然肩并肩
（相互接触！）但现在可以来回穿插、彼此滑过。他们仍须
与人群保持接触——谁也不许自由跑开。
\textbf{3.} 学生四散奔跑，快速地沿直线前进，直到
相撞或撞到墙壁，然后不停地反弹——既然没有东西
把他们连在一起，只有墙壁能让他们停下，人群
占满了体育馆的每个角落。
\textbf{4.} 是的——每一幕都是同样的三十人：\omterm{def:g3:measuring-mass:mass}{质量}
在物态变化中保持不变，这是舞台上演示的。
\textbf{5.} \omterm{def:g3:solids-liquids-gases:gas}{气体}一幕允许这样做——跑步者挤进
体育馆的一半（并且把牧羊人撞得更猛）。\omterm{def:g3:solids-liquids-gases:liquid}{液体}那一团不能缩小：
学生已经相互接触。演示：\omterm{def:g3:solids-liquids-gases:gas}{气体}可被\omterm{prop:g7:states-of-matter:squeeze}{压缩}，
\omterm{def:g3:solids-liquids-gases:liquid}{液体}则不能。
\textbf{6.} 在各自位置上越抖越剧烈；
抖动冲破队列的那一刻，这一幕便成了\omterm{def:g3:solids-liquids-gases:liquid}{液体}，
学生开始彼此滑过。
\textbf{7.} 刻度盘的\omterm{def:g5:everyday-energy:energy}{能量}用于打破堆叠的束缚——
一排一排地打破——而不用于加快运动：\omterm{def:g3:temperature-thermometers:temperature}{温度计}评判员
必须报告，在最后一排瓦解之前，平均活跃程度\emph{保持不变}。
平台期，就此演示出来。
\textbf{8.} 只有\emph{最快}的学生——只有他们
能够挣脱。其余人群的平均活跃程度下降：\omterm{def:g3:solids-liquids-gases:liquid}{液体}
通过\omterm{def:g2:ice-water-steam:evapcond}{蒸发}使自己变凉，
与汗水和湿手指的凉意如出一辙。
\textbf{9.} 演员不断与绳子相撞，并将它向外推；
\omterm{def:g5:speed-distance-time:speed}{速度}更快的演员推得更猛、更频繁，圆环便膨胀起来：
正如阳光下的气球。
\textbf{10.} “以\emph{疏松}的次序堆叠——手臂的距离，
空隙多”：水的\omterm{def:g3:solids-liquids-gases:solid}{固体}比\omterm{def:g3:solids-liquids-gases:liquid}{液体}占据更多空间，因此
冰一幕中的木筏会浮在\omterm{def:g3:solids-liquids-gases:liquid}{液体}一幕的人群之上。
\textbf{11.} 蓝色衬衫被推挤——一步步随机地
撞过滑动的人群，直到游荡遍各处：墨水在扩散。
把刻度盘调高，每次推挤都更快，
游荡也就更快。
\textbf{12.} 例如：“今晚的演出只呈现一个观点：一切物质
都是运动中的\omterm{def:g7:states-of-matter:molecule}{分子}。诚然，两条旧定律也随之而来——
\omterm{def:g3:measuring-mass:mass}{质量}在每一次物态变化中保持不变，以及\omterm{def:g6:water-states:changes}{熔化}时的
平台期。我们无法演出的是尺度：真实的演员
比任何能被看见的东西都小，
数量比世界上所有的观众还多。”
\end{solution}
